\section{The shared-point quartic in 1272 variables}
\label{sec:v5-1272}

This section gives a new complete quadratic compilation of the shared-point
construction in Jos\'e Brox's public manuscript~\cite{Brox2026}.
The shared comparison points, reciprocal-trace detectors, fixed height
scales, direct Lucas test, and affine anchor change are Brox's arithmetic
and algebraic improvements. The simultaneous prime theorem used to choose
the shared points is Sun's Theorem~4.1, in the precise form of
Theorem~\ref{thm:v5-simultaneous-primes} below. Our contribution here is an
independently specified gate program, its source verification, and the
complete polynomial obtained after 32 certified eliminations. The
definition does not depend on an unavailable coefficient table.
The public source passages are Brox's
\href{https://github.com/JoseBrox/diophantine-z-q_improvements/blob/f4d1b5e6c23d4adeb3ae1f986ab0bf87884bdff6/paper/diophantine_optimization_20261106_v3.tex#L2354-L3044}{shared-point construction},
\href{https://github.com/JoseBrox/diophantine-z-q_improvements/blob/f4d1b5e6c23d4adeb3ae1f986ab0bf87884bdff6/paper/diophantine_optimization_20261106_v3.tex#L3352-L3405}{fixed height scales},
and
\href{https://github.com/JoseBrox/diophantine-z-q_improvements/blob/f4d1b5e6c23d4adeb3ae1f986ab0bf87884bdff6/paper/diophantine_optimization_20261106_v3.tex#L4135-L4758}{quadratic compilation and anchor translation}.

We retain the arithmetic conclusions and conventions of
Section~\ref{sec:degree8-original}, including the restricted-converse
input to Lemma~\ref{hyp:o8-pole}. In this section the height estimate is
Theorem~\ref{hyp:o8-height}, with exponent 24 and its original constant
$H_*$; it is not the exponent-19 estimate used for the preceding quartic.
Thus
\[
  1\le H_*<2^{4783227}.
\]
The exact finite compilation and coefficient assertions below require no
arithmetic conjecture. Identifying the rational input projection with
$\Z$ uses the stated arithmetic results and the simultaneous prime theorem.

\begin{theorem}[The new quartic]\label{thm:v5-1272-polynomial}
The complete table specified below defines a polynomial
\[
 P_{1272}\in\Z[X_0,\ldots,X_{1271}]
\]
of degree four, with exactly $48001$ nonzero coordinate monomials,
integer coefficient content one, and coefficient height
\[
 H(P_{1272})=64\cdot3^{2^{28}}<2^{429496736}.
\]
It is the sum of the squares of 1229 quadratic residuals. Under the
arithmetic conclusions and conventions just stated, for every $X_0\in\Q$,
\[
 X_0\in\Z\quad\Longleftrightarrow\quad
 \exists X_1,\ldots,X_{1271}\in\Q:\ P_{1272}(X_0,\ldots,X_{1271})=0.
\]
In particular its 1272 variables include the input and 1271 existential
witnesses.
\end{theorem}

\subsection{Points, local parameters, and quotient domains}
\label{subsec:v5-point-source}

Use the normalized curve and generator already established in
Section~\ref{sec:degree8-original}:
\[
 K=\Q(\sqrt2),\quad d=3-\sqrt2,\quad D=d^2,
 \quad E:y^2=x^3-Dx,
\]
\[
 P_*=(57/8-7\sqrt2/2,-165/8+417\sqrt2/32),
 \qquad 2E(K)=\Z P_*.
\]
The last assertion follows from Proposition~\ref{prop:generator} by the
unit change of variables recorded there. Conjugation of $K/\Q$ is
denoted by $\sigma$. All field equations below mean both rational
component equations in the basis $1,\sqrt2$.

The source uses nine nonidentity points
\[
 A,S,T,Z,H,U_A,U_S,U_T,U_Z\in2E(K).
\]
Each is represented as the double of an affine halfpoint
$(\alpha,\eta)\in E(K)$, using four rational halfpoint coordinates.
The equations for a double are
\begin{equation}\label{eq:v5-double}
 \eta^2=\alpha^3-D\alpha,\qquad
 2\eta m=3\alpha^2-D,\qquad
 x=m^2-2\alpha,\qquad y=m(\alpha-x)-\eta.
\end{equation}
Products of degree greater than two in this display are evaluated by
the quadratic gates specified below. Set
\[
 Q_0=P_*,\quad(Q_1,Q_2,Q_3,Q_4)=(A,S,T,Z),
 \quad(U_0,U_1,U_2,U_3,U_4)=(H,U_A,U_S,U_T,U_Z).
\]
Compute $[4]Q_i$ by two doublings, or by fixed rational arithmetic for
$Q_0$, and compute
\[
 V_i=U_i-[4]Q_i\quad(0\le i\le4).
\]
For the subtraction of $(x',y')$ from $(x,y)$ we retain an inverse of
$x-x'$ in $K$ and use
\[
 m=(y+y')/(x-x'),\quad x''=m^2-x-x',\quad
 y''=m(x-x'')-y.
\]

\begin{lemma}[Domains of the point quotients]\label{lem:v5-point-domains}
The equations \eqref{eq:v5-double} force $\eta\ne0$ and describe a
finite nonidentity point of $2E(K)$. Every such point has nonzero
abscissa and ordinate. The guarded differences $V_i$ have the same
properties. Consequently $f(Q)=1/x(Q)$ and $z(Q)=-x(Q)/y(Q)$ are
defined and nonzero at every point where they are used here.
\end{lemma}
\begin{proof}
If $\eta=0$, the curve equation gives $\alpha=0$ or $\alpha^2=D$.
The right side of the slope equation is then $-D$ or $2D$, respectively,
a contradiction. The ordinary duplication formulas therefore apply.
They give
\[
 x([2](\alpha,\eta))=
 \left(\frac{\alpha^2+D}{2\eta}\right)^2\ne0,
\]
because $D=d^2$ and the numerator is positive at either real embedding.
An ordinate-zero point has abscissa $0,d$, or $-d$.
The latter two are not squares in $K$: their rational norm is 7,
whereas the norm of a square is a rational square. Thus the doubled
ordinate is nonzero. A guarded affine difference is finite and is
again in $2E(K)$, so the same argument applies. Conversely any
nonidentity point of $2E(K)$ has an affine halfpoint of nonzero
ordinate and is represented by these formulas.
\end{proof}

Write
\[
 a_0=1,\quad a_1=-2/5,\quad
 a_j=\frac{4j-5}{4j+1}a_{j-1}\ (j\ge2),\qquad
 G_J(w)=\sum_{j=0}^J a_jw^j,
\]
and, for $e=16,52$, put
\begin{equation}\label{eq:v5-log}
 \ell_e(Q)=z(Q)G_{e/4-1}(Df(Q)^2).
\end{equation}
Lemma~\ref{lem:quartic-reciprocal-nonzero} applies to both truncation
degrees 3 and 12. Thus all these values are nonzero globally.

\begin{lemma}[The logarithm estimate at unramified places]
\label{lem:v5-unramified-log}
Let $p>76176$ be a rational prime and $\mathfrak p\mid p$ a place
of $K$. If a nonidentity point $Q$ is in the formal group at
$\mathfrak p$, normalize $v(p)=1$ and put $r=v(z(Q))$. Then $r\ge1$,
and for $e\in\{16,52,64\}$,
\[
 v(\log_E Q)=r,\qquad
 v(\log_E Q-\ell_e(Q))\ge(e+1)r,\qquad
 \ell_e(Q)/\log_E Q\in1+p^e\mathcal O_{K_{\mathfrak p}}.
\]
Here $\ell_{64}$ is defined by the same formula with $G_{15}$.
Ratios with numerator precision $e$ and denominator precision $e'$
in this set have the valuation of their full logarithm ratio and
relative error of valuation at least $\min(e,e')$.
\end{lemma}
\begin{proof}
The local extension is unramified, including when $p$ is inert.
The curve has good reduction and $D$ is a unit. On the formal branch,
\[
 f=z^2(1-Df^2),\qquad w=Df^2,\qquad v(w)=4r.
\]
The recurrence gives
$4w(1-w)G'(w)+(1+w)G(w)=1-w$ and hence the same identity
$\log_E Q=zG(w)$ as in Lemma~\ref{lem:quartic-reciprocal-log}.
For $j\ge1$ the denominator of $a_j$ divides $(4j+1)!$, so
\[
 v(a_j)\ge-\frac{4j+1}{p-1}.
\]
Set $m=e/4\le16$. If $j\ge m$ and $4j+1<p$, then
$v(a_jw^j)\ge4mr$. Otherwise $j\ge19044$ and
\[
 v(a_jw^j)\ge4jr-j\ge3jr\ge4mr.
\]
The lower bounds tend to infinity, proving the estimate for the
entire tail. Every nonconstant term of $G(w)$ has positive valuation,
so the leading term gives $v(\log_E Q)=r$. Division proves the
relative estimate, and division of two relative estimates proves
the ratio assertion. No step uses residue degree one.
\end{proof}

\subsection{The complete rational source}
\label{subsec:v5-complete-source}

The ordinary input is $X$, and $\beta=X^2+2$. Define
\[
 (k_1,k_2,k_3,k_4)=(16,52,52,52),\qquad
 L_i=\frac{\ell_{k_i}(U_i)}{\ell_{52}(H)},\qquad
 L'_i=\sigma\!\left(\frac{\ell_{k_i}(V_i)}{\ell_{52}(V_0)}\right).
\]
The sixteen field tests $D_j$, with their orders $m_j$, are exactly
\begin{equation}\label{eq:v5-test-list}
\begin{array}{c|c}
 f(U_0),\ldots,f(U_4),\ \sigma f(V_0),\ldots,\sigma f(V_4)&1\quad(10)\\
 L_i-L'_i\ (1\le i\le4)&16,52,52,52\\
 L_1-\beta&16\\
 F(L_2,L_3,L_4)&52,
\end{array}
\end{equation}
where
\begin{equation}\label{eq:v5-lucas-polynomial}
 F(U,V,W)=V^2-U(U-2)VW+(U-2)^2W^2-1.
\end{equation}
There is no rational proxy for an integer point label.

For $j=1,\ldots,4$ write $x(Q_j)=x_{j,0}+x_{j,1}\sqrt2$ and set
\[
 \mathcal A_0=1,\quad
 \mathcal A_{2j-1}=x_{j,0}^2,\quad
 \mathcal A_{2j}=(1+x_{j,1}^2)^2.
\]
These are positive nonzero rational squares. In fact $x(Q_j)$ is
positive at both real embeddings by Lemma~\ref{lem:v5-point-domains},
so $x_{j,0}>0$. Retain one rational coordinate $\lambda$ and put
\[
 \Theta=\{0,-1/2,-1,1/2399,-32/407\},\qquad
 F_5(T)=\prod_{\theta\in\Theta}(T-\theta),\qquad
 t_r=\mathcal A_r\lambda,\quad\delta_r=t_r-1.
\]
For every $0\le r\le8$, require $t_r\notin\Theta\cup\{1\}$ and define
\begin{align}
 R_r&=\mathcal A_r^{-5}F_5(t_r)/\delta_r^6,
 &\widetilde D_{r,j}&=D_j/\delta_r^{m_j+1},\label{eq:v5-detector-data}\\
 W_r&=\prod_{j=1}^{16}
 \left(1+R_r(\widetilde D_{r,j}/R_r^{m_j})^2\right),
 &T_r&=\Tr_{K/\Q}(W_r^{-1}),\qquad B_r=R_r/T_r^2.
 \label{eq:v5-detector-product}
\end{align}
Require $W_r,T_r\ne0$. For $B=B_r$ retain four rational witnesses
$e_r,s_r,p_r,q_r$ and impose
\begin{equation}\label{eq:v5-zero-pole}
 \begin{aligned}
 C_r&=17B_r,\qquad h_r=e_r(C_r-B_r^2),\\
 p_r^2&=C_r(1-e_r)(1+3e_r),\\
 C_rq_r^2&=(s_r^2-C_rh_r)(s_r^2+3C_rh_r).
 \end{aligned}
\end{equation}
The first two expressions are arithmetic definitions. The last two
are the preparation-zero specialization of the original pole predicate.
There is no nonzero guard on $e_r,s_r,p_r,q_r$, or $h_r$.

The scalar block shares the point $H$. Define
\[
 r_\beta=\beta(\beta-1)(\beta-2401)(32\beta+343)(2\beta-1).
\]
It is nonzero for rational $X$: the only positive possible root not
excluded by $\beta\ge2$ would require $X^2=2399$.
Let $Y_1,\ldots,Y_6$ be, in order, the two rational components of
$f(H),f(S),z(S)/z(H)$, and set
\[
 (e_1,\ldots,e_6)=(1,1,1,1,32,32),\qquad
 y_i=Y_i/r_\beta^{e_i},\qquad N_r(u,v)=u^2+rv^2.
\]
The scalar norm tree has no padding leaf. It is
\begin{equation}\label{eq:v5-scalar-tree}
 \begin{aligned}
 n_0&=N_{r_\beta}(1,y_1),\quad n_1=N_{r_\beta}(y_2,y_3),\quad
 n_3=N_{r_\beta}(y_5,y_6),\\
 W_\beta&=N_{r_\beta}\bigl(N_{r_\beta}(n_0,n_1),
                          N_{r_\beta}(y_4,n_3)\bigr).
 \end{aligned}
\end{equation}
Use eight Boolean coordinates $\epsilon_-,\epsilon_p$, for
$p\in\{2,3,5,7,11,13,17\}$, impose $\epsilon^2-\epsilon=0$, and set
\[
 c=(1-2\epsilon_-)\prod_p(1+(p-1)\epsilon_p),\qquad
 b_\beta=c r_\beta W_\beta^2,\qquad C_\beta=17b_\beta.
\]
Require $b_\beta\ne0$. Retain five rational witnesses
$t_\beta,e_\beta,s_\beta,p_\beta,q_\beta$ and impose
\begin{align}
 (1+b_\beta t_\beta^2)U_\beta
    &=b_\beta(1-b_\beta t_\beta^2),
 &h_\beta&=e_\beta(C_\beta-U_\beta^2),\label{eq:v5-scalar-preparation}\\
 p_\beta^2&=C_\beta(1-e_\beta)(1+3e_\beta),
 &C_\beta q_\beta^2&=(s_\beta^2-C_\beta h_\beta)
                     (s_\beta^2+3C_\beta h_\beta).
 \label{eq:v5-scalar-pole}
\end{align}
Here $U_\beta,h_\beta$ are arithmetic values. The preparation equation
forces $1+b_\beta t_\beta^2\ne0$: otherwise $b_\beta t_\beta^2=-1$
but its right side is $2b_\beta\ne0$.

Finally fix the integers
\begin{equation}\label{eq:v5-fixed-constant}
 C=3^{2^{26}},\qquad \mathsf S=C^2=3^{2^{27}},\qquad
 \kappa=\mathsf S^2=3^{2^{28}}.
\end{equation}
The symbol \texttt{S} in the coefficient files means $\mathsf S$;
it is not the point $S$ and is not a quantified coordinate.
Write $f(S)=f_0+f_1\sqrt2$, retain three private rational witnesses
$u,v,w$, and impose
\begin{equation}\label{eq:v5-physical-anchor}
 u^2+6v^2-2w^2-3(f_0^2+f_0f_1-f_1^2)^2+2\mathsf S^2=0.
\end{equation}
This completes the rational source. Every rational expression has the
displayed guard or a domain implication already proved. Equations
\eqref{eq:v5-zero-pole} and \eqref{eq:v5-scalar-pole} retain their
degenerate witness values.

\begin{lemma}[The three-witness anchor]\label{lem:v5-three-witness-anchor}
Equation~\eqref{eq:v5-physical-anchor} is equivalent to
$f_0/C,f_1/C\in\Z_{(2)}\cap\Z_{(3)}$. For the nonzero integer label
$a$ of $S=[a]P_*$ it implies
\[
 0\ne a\in2\,3^{2^{25}-1}\Z,\qquad
 |a|>(21H_*)^{13/2}>(20H_*)^4.
\]
Its permitted labels contain arbitrarily divisible positive integers.
\end{lemma}
\begin{proof}
Put $R=\Z_{(2)}\cap\Z_{(3)}$. The quaternion algebra $(2,3)_\Q$
ramifies exactly at 2 and 3. The quadratic embedding criterion
\cite[Proposition 14.6.7]{voight2021} shows, for rational $t\ne0$,
that $1-2/(3t^2)$ is represented by $2x^2+3y^2-6z^2$ precisely
when it is nonsquare in both $\Q_2$ and $\Q_3$. This last condition
is exactly $t\in R$. Indeed at 2 a positive valuation of $t$ gives
odd valuation of the expression, valuation zero gives residue 3
modulo 8, and negative valuation gives a square in $1+8\Z_2$.
At 3, nonnegative valuation of $t$ gives odd negative valuation,
while negative valuation gives a square in $1+3\Z_3$.
The expression is nonzero because $t^2=2/3$ has no rational solution.
Multiplying by $3t^2$ gives
\[
 t\in R\quad\Longleftrightarrow\quad
 \exists u,v,w\in\Q:\ u^2+6v^2-2w^2-3t^2+2=0.
\]
The case $t=0$ has witnesses $(0,0,1)$. A rational-square value of
$1-2/(3t^2)$ is excluded by the division-algebra criterion as well:
a nonzero trace-zero quaternion cannot square to a nonzero rational
square in a division algebra.

The form $A^2+AB-B^2$ has no nontrivial zero over $\F_2$ or $\F_3$.
Scaling a pair therefore proves
$v_p(a^2+ab-b^2)=2\min(v_p(a),v_p(b))$ for $p=2,3$.
Apply the preceding three-witness formula to
$(f_0^2+f_0f_1-f_1^2)/C^2$ and multiply its witnesses by $C^2$.
This gives exactly \eqref{eq:v5-physical-anchor}.

The local calculation in Section~\ref{subsec:o8-anchor-subgroup}
applies with exponent $2^{26}$: the reduction of $P_*$ at 3 has
order two and $v_3(z(2P_*))=1$. For even nonzero $a$ it gives
$v_3(f([a]P_*))=2+2v_3(a)$. At 2 the rational components of
$f([a]P_*)$ are integral. The anchor is consequently equivalent to
$0\ne a\in2\,3^{2^{25}-1}\Z$. Its minimum possible absolute value
has binary logarithm greater than
\[
 1+\tfrac32(2^{25}-1)=50331647.5,
\]
whereas
$\log_2((21H_*)^{13/2})<\tfrac{13}{2}(5+4783227)=31091008$.
The last comparison with $(20H_*)^4$ follows from
$21^{13}>20^8$ and $H_*\ge1$. The subgroup description proves the
last assertion.
\end{proof}

\subsection{Local tests and the height implication}
\label{subsec:v5-soundness}

\begin{lemma}[One place for every test]\label{lem:v5-one-place}
Let $R\in\Q^*$, $D_j\in K$, and positive integers $m_j$ define
\[
 W=\prod_j(1+R(D_j/R^{m_j})^2),\quad
 T=\Tr(W^{-1}),\quad B=R/T^2,
\]
with $W,T\ne0$ and a preparation-zero pole certificate for $B$.
If $p>76176$ and $v_p(R)$ is odd, then $v_p(R)>0$ and some
$\mathfrak p\mid p$ has $v_{\mathfrak p}(D_j)\ge m_j$ for every $j$.
\end{lemma}
\begin{proof}
Set $\alpha=v_p(R)$. The nonconstant term in each nontrivial factor
has odd valuation $\alpha+2(v(D_j)-m_j\alpha)$; unramifiedness
makes these valuations integral. The factor valuation is its minimum
with zero. Thus $\lambda_{\mathfrak p}=v_{\mathfrak p}(W)\le0$.
Put $\lambda=\max_{\mathfrak p\mid p}\lambda_{\mathfrak p}$.
Integrality of trace gives $v_p(T)\ge-\lambda$. Hence
\[
 v_p(B)\le\alpha+2\lambda.
\]
The left side is odd and, by Lemma~\ref{hyp:o8-pole}, not negative.
It is positive, so $\alpha>0$ and $\lambda>-\alpha/2$.
At a place attaining $\lambda$, each factor has valuation at least
$\lambda$, because all factor valuations are nonpositive. For
nonzero $D_j$ this implies
\[
 v(D_j)>(m_j-3/4)\alpha\ge m_j-3/4.
\]
Integrality gives $v(D_j)\ge m_j$; a zero $D_j$ satisfies it as well.
\end{proof}

\begin{lemma}[The two-anchor comparison]\label{lem:v5-two-anchors}
Write the integer labels of $(A,S,T,Z)$ as $(n,a,t,b)$.
If the first fourteen tests in \eqref{eq:v5-test-list} have their
specified orders at a place $\mathfrak p\mid p$, $p>76176$, then
\[
 L_i\in\mathcal O_{K_{\mathfrak p}},\qquad
 L_i\equiv n_i\pmod{p^{k_i}},\qquad
 (n_1,n_2,n_3,n_4)=(n,a,t,b).
\]
This holds at split and inert primes.
\end{lemma}
\begin{proof}
Write $H=[h]P_*$ and $U_i=[u_i]P_*$. The order-one tests put
$H,U_i$ in the formal group at $\mathfrak p$ and $V_0,V_i$ in the
formal group at the conjugate place. Positive valuation of $f=1/x$
gives formal membership on this good model. The guarded points give
$h\ne0,4$. Their full logarithm ratios, after conjugating the second,
are
\[
 r_i=u_i/h,\qquad r'_i=(u_i-4n_i)/(h-4).
\]
These identities follow from the homomorphism property on equal
formal multiples, and do not require $P_*$ itself to be formal.
At least one of $h,h-4$ is a $p$-unit. Its ratio is integral, and
Lemma~\ref{lem:v5-unramified-log} gives the same valuation to its
truncated ratio. The cross test makes the other truncated ratio
integral too. Both relative errors are therefore absolute errors,
and $L_i\equiv r_i\equiv r'_i$ modulo $p^{k_i}$. Multiplication
by $h,h-4$ and subtraction give $4L_i\equiv4n_i$. Divide by four.
At an inert prime conjugation preserves the single local valuation,
so the same argument applies.
\end{proof}

\begin{lemma}[Height transfer]\label{lem:v5-height-transfer}
Let $M_r$ be the odd-contact radical of $t_r$ for the five contacts
$\Theta$, omitting primes at most 76176. If $M_r^k$ divides a
nonzero integer $N$ for every $r$, then
\[
 h(x(Q_i))<7H_*|N|^{24/k}\qquad(1\le i\le4).
\]
At every prime counted in $M_r$, the source imposes all sixteen
field tests with their specified orders at one place above that prime.
\end{lemma}
\begin{proof}
The fractional linear map $\phi(t)=2+1/t$ carries $\Theta$ to
$\infty,0,1,2401,-343/32$. Its integral matrix has determinant $-1$.
All primes dividing the relevant fixed coefficients and contact
differences, including differences from 1, lie in
$\{2,3,5,7,11,37,109,439,2399\}$. Thus outside the cutoff
$M_r=M_5(\phi(t_r))$, and
\[
 h(t_r)\le h(\phi(t_r))+\log3
       \le(H_*+\log3)M_r^{24}.
\]
At a prime of $M_r$, $\delta_r$ is a unit. Exactly one contact
factor has odd positive valuation, while $\mathcal A_r^{-5}$ is
a rational square. Therefore $v_p(R_r)$ is odd.
Lemma~\ref{lem:v5-one-place} applied to $\widetilde D_{r,j}$ gives
the asserted orders for $D_j$.

Set $T_0=(H_*+\log3)|N|^{24/k}$. The case $r=0$ gives
$h(\lambda)\le T_0$, and the other cases give
$h(\mathcal A_r\lambda)\le T_0$. Hence
$h(\mathcal A_r)\le2T_0$. For the first component,
$h(x_{i,0})\le T_0$. For $x_{i,1}=a/b$ in lowest terms,
\[
 h(1+x_{i,1}^2)=\log(a^2+b^2)\ge2h(x_{i,1}),
\]
so $h(x_{i,1})\le T_0/2$. Finally
\[
 h(x(Q_i))\le h(x_{i,0})+h(x_{i,1})+\tfrac32\log2
 <7H_*|N|^{24/k}.
\]
\end{proof}

\begin{lemma}[The two reduced group cardinalities]
\label{lem:v5-two-reductions}
Let $p>76176$ be prime and let $\pi\in\Z[i]$ have norm $p$, with
$\pi\equiv9\pmod{16}$ and $\pi\equiv1+3i\pmod7$.
The two reductions of $E/K$ at the places above $p$ have group
cardinalities $N(\pi-1)$ and $N(\pi+1)$, in some order.
\end{lemma}
\begin{proof}
The congruences give $p\equiv1\pmod8$ and $p\equiv3\pmod7$.
Thus $p$ splits in $K$ and $(7/p)=-1$. The reductions of
$d=3-\sqrt2$ have product 7 and opposite square classes. The two
curves $y^2=x^3-d^2x$ are consequently opposite quadratic twists
of $y^2=x^3-x$ and have opposite traces.

The Frobenius calculation already given in
Section~\ref{subsec:o8-pattern-proof} applies: the order-four
automorphism is defined over $\F_p$, its rational centralizer is
$\Q(i)$, and integrality puts Frobenius in $\Z[i]$.
Write it as $a+bi$, so $a^2+b^2=p$ and the two cardinalities are
$p+1-2a$ and $p+1+2a$. Full rational two-torsion makes both
cardinalities divisible by four. Since $p\equiv1\pmod4$, this
forces $a$ odd and $b$ even. The real part of $\pi$ is odd and its
imaginary part is even. Unique factorization in $\Z[i]$ therefore
makes $a+bi$ one of $\pm\pi,\pm\overline\pi$; the associates
with factor $i$ have the wrong parity. Hence the unordered
cardinalities are exactly $N(\pi-1),N(\pi+1)$.
\end{proof}

\begin{theorem}[Arithmetic interpretation of the source]
\label{thm:v5-source-projection}
With the inherited arithmetic results and
Theorem~\ref{thm:v5-simultaneous-primes}, the rational source of
Section~\ref{subsec:v5-complete-source} has input projection exactly $\Z$.
\end{theorem}
\begin{proof}[Soundness]
Let $\mathcal B=\max(|n|,|a|,|t|,|b|)$ and
$\Delta=F(a,t,b)\in\Z$. Polynomial substitution preserves
congruences between integral local elements. Lemmas
\ref{lem:v5-two-anchors} and \ref{lem:v5-height-transfer} therefore
give $M_r^{52}\mid\Delta$ whenever $\Delta\ne0$.
The index-height estimate \eqref{eq:o8-indexheight} yields
\[
 \mathcal B^2-6<7H_*|\Delta|^{24/52},\qquad
 |\Delta|<5\mathcal B^4,
\]
the latter because $\mathcal B\ge|a|>7$. Consequently
\[
 \mathcal B^{2/13}<(7\cdot5^{6/13}+6)H_*<21H_*.
\]
For the last strict inequality, $5^6 7^{13}<15^{13}$.
This contradicts Lemma~\ref{lem:v5-three-witness-anchor}. Thus
$\Delta=0$, and Lemma~\ref{lem:o8-lucas}, with $b\ne0$, gives
$\log\mathcal B>|a|-3$.

Lemma~\ref{v5:scalar32-bridge} applies to the exact tree
\eqref{eq:v5-scalar-tree} and gives
\[
 h(\beta)\le H_*|a|^{3/4}<|a|/20.
\]
If $\beta\ne n$, put $N=|\operatorname{num}(\beta-n)|$ in lowest
terms. The order-16 test gives $M_r^{16}\mid N$. Therefore
\[
 \mathcal B^2<7H_*N^{3/2}+6,\qquad
 N\le2e^{h(\beta)}\mathcal B,
\]
and $\log\mathcal B<2\log(32H_*)+3h(\beta)$.
The anchor also gives $\log(32H_*)<|a|/20$, so the right side
is less than $|a|/4<|a|-3$, a contradiction. Hence
$\beta=n\in\Z$. Since $X\in\Q$ and $X^2=n-2\in\Z$,
the reduced denominator of $X$ is one.
\end{proof}

\begin{proof}[Completeness]
Fix $X\in\Z$, set $n=X^2+2$, and choose the Boolean correction
$c$ so that $cr_\beta>0$ is square in $\Q_2,\Q_3,\Q_{17}$,
as in the finite squareclass computation accompanying
Lemma~\ref{hyp:o8-pole}. Let $\mathcal S$ be the prime support of
$cr_\beta$. The scalar completeness construction in
Lemma~\ref{v5:scalar32-bridge} first chooses $h_0>0$ divisible by
64 with $[h_0]P_*$ sufficiently deep at these primes. It then
chooses a positive label $a$ sufficiently divisible relative to
$h_0$ and satisfying the anchor. The normalized scalar leaves are
integral at primes dividing $r_\beta$ and divisible at primes
dividing $c$ but not $r_\beta$. Choose positive $t,b$ with
$F(a,t,b)=0$ by Lemma~\ref{lem:o8-lucas}.
This fixes the four root points and all scales $\mathcal A_r$.

Temporarily put $H=[h]P_*$ and $U_i=[hn_i]P_*$.
All comparison tests are rational functions of $[h]P_*$ with
nonzero denominators at $h=h_0$. Their poles lie above finite
sets under fixed nonzero multiplication and translation maps
on the two real factors, as required by the real progression lemma.
Continuity preserves the scalar
conditions at every place over $\mathcal S$ and bounds below the
valuations of all sixteen comparison tests at the places over
$2,3,17$, throughout a progression
\[
 h\equiv h_0\pmod L,\qquad64\mid L.
\]
One obtains $L$ by putting sufficiently small open formal subgroups
inside the finitely many translated neighborhoods and taking a
common multiple of their indices. Lemma~\ref{v5:real-progression-bridge}
gives, independently of future progression moduli, real bounds for
the tests attained by arbitrarily large positive $h$ in each
nonconstant progression. The direct polynomial $F$ introduces
no denominator and has the same boundedness property.

Choose a positive integer $D_0$ such that every
$D_0^2\theta/\mathcal A_r$ is integral, and put
\[
 v=D_0^2w^2,\qquad\lambda=u/v.
\]
The contact positions are the integer forms
$u-(D_0^2\theta/\mathcal A_r)w^2$. Delete repeated forms, but retain
their incidences in the nine detectors. The five forms within one
detector are distinct, and the zero contact is the common form $u$.
Apply Theorem~\ref{thm:v5-simultaneous-primes} to these forms.
Prescribe arbitrarily high divisibility of $w$, a sufficiently
large lower bound for $u/w^2$, and primes larger than 76176 and
all prime factors of the fixed scales and contact differences.
In its odd squarefree avoidance modulus include every odd prime
factor of $3\cdot17L$ except 7.
Writing this modulus as $M$, the prescribed Gaussian classes are
\[
 \pi_p\equiv9\pmod{16},\qquad
 \pi_p\equiv1+3i\pmod7,\qquad
 \pi_p\equiv i\pmod M.
\]

Let the resulting Gaussian generators be $\pi_p$, and write
$d_-(p)=N(\pi_p-1)$ and $d_+(p)=N(\pi_p+1)$ for the cardinalities
of the two reduced elliptic groups, orienting the places by
Lemma~\ref{lem:v5-two-reductions}. These cardinalities annihilate
the reductions of $P_*$; their exact point orders are not needed.
They satisfy
\[
 v_2(d_-)=6,\quad v_2(d_+)=2,\quad
 \gcd(d_-(p),d_+(q))=4
\]
for every pair of selected primes. No odd prime factor of $L$
divides either order. For 7 this follows separately from
$d_+\equiv6$ and $d_-\equiv2\pmod7$.
Set $M_-=\operatorname{lcm}_p d_-(p)$ and
$M_+=\operatorname{lcm}_p d_+(p)$. The congruences
\[
 h\equiv0\pmod{M_-},\quad h\equiv4\pmod{M_+},\quad
 h\equiv h_0\pmod L
\]
are compatible: their pairwise gcds are $4,64,4$, respectively.
Choose a sufficiently large positive $h$ in this progression
attaining the real bounds. It avoids $0,4,-4$, so all point and
subtraction guards hold. At a selected prime, $H,U_i$ are formal
at the place of order $d_-$, while $V_0,V_i$ are formal at its
conjugate place. Their exact ratios are $n_i$. Lemma
\ref{lem:v5-unramified-log} gives all sixteen required orders,
including the order-52 polynomial test because $F(a,t,b)=0$.

We check that the nine preparation-zero poles can be filled.
For each detector,
\[
 R_r=\frac{\prod_{\theta\in\Theta}
                   (u-(\theta/\mathcal A_r)v)}
                 {v^5\delta_r^6}.
\]
The denominator is a rational square. At the five distinct contact
primes of this detector $R_r$ has valuation one, and elsewhere
its valuation is even. The primes do not divide $v$: otherwise
two distinct forms in the same detector would give the same
prime value, impossible since $w>0$ and their coefficients are
distinct. The lower bound for the primes excludes the fixed scale
and difference denominators. At either place over a
contact prime a factor of $W_r$ has valuation zero or negative,
and its inverse has residue one or zero, respectively. At the
successful place every inverse factor has residue one.
Consequently $T_r$ has residue one or two and is a unit.
It follows that $B_r$ has positive odd valuation at its contact
primes and even valuation elsewhere. Its reduced denominator is
a square.

The prescribed Gaussian classes make every contact prime square
in $\Q_2,\Q_3,\Q_{17}$. The last product formula therefore makes
$B_r$ square in each of those fields. The uniform real bounds
and $\lambda\longrightarrow+\infty$ give
\[
 R_r\sim\mathcal A_r^{-6}\lambda^{-1},\quad
 W_r=1+O(\lambda^{-3}),\quad T_r=2+O(\lambda^{-3}),\quad
 B_r\sim(4\mathcal A_r^6\lambda)^{-1}.
\]
In particular $0<B_r<17$. At a small prime $p=2,3,17$, put
$\nu_p=v_p(v)$ and $\eta_{p,r}=v_p(\mathcal A_r)$.
The common form $u$ is a selected prime greater than these small
primes, so $v_p(u)=0$ and $v_p(\lambda)=-\nu_p$.
The fixed neighborhoods give $v(D_j)\ge-c_p$. For sufficiently
large prescribed $\nu_p$ one has
\[
 v_p(R_r)=\nu_p-6\eta_{p,r},\qquad
 v(\widetilde D_{r,j})\ge-c_p+(m_j+1)(\nu_p-\eta_{p,r}).
\]
The nonconstant factor term has valuation at least
$3\nu_p+(10m_j-8)\eta_{p,r}-2c_p$, tending to infinity.
Thus $T_r$ is as close to 2 as needed, including at the ramified
prime 2, and
\[
 v_p(B_r)=\nu_p-6\eta_{p,r}-2v_p(2).
\]
Prescribing the divisibility of $w$ before applying the prime theorem
therefore ensures $v_2(B_r)\ge3$, $v_3(B_r)\ge1$,
$v_{17}(B_r)\ge2$. Lemma~\ref{v5:preparation-zero-bridge} supplies
the four witnesses in every coordinate pole.

Finally the progression preserved all scalar leaf conditions.
The constant leaf on the all-left path makes $W_\beta$ a unit at
every prime dividing $r_\beta$; the divisibility of the other
leaves makes it a unit at primes dividing $c$ but not $r_\beta$.
Since $|r_\beta|>1$, it is nonzero. Therefore
$b_\beta=cr_\beta W_\beta^2$ is positive, has square reduced
denominator, and is locally square at $2,3,17$.
Lemma~\ref{hyp:o8-pole} supplies the scalar preparation and pole
witnesses for this final $H$. The anchor already holds for $S$.
All witnesses have now been supplied without changing the root points.
\end{proof}

\subsection{Exact quadratic compilation and the complete table}
\label{subsec:v5-compilation}

The implementation retains an affine rational expression whenever
possible. For a quadratic expression it writes
\[
 q=cq_2+q_1,
\]
where $q_1$ is affine and $q_2$ is the homogeneous quadratic part
normalized by its first nonzero coefficient. A coordinate $v$
with equation $v-q_2=0$ is introduced once for each distinct
normalized $q_2$, and the value $cv+q_1$ is reused. This exact
sharing rule includes additions and fixed rational scalings without
new coordinates. A field product is assigned through
\[
 v_0=a_0b_0+2a_1b_1,\qquad v_1=a_0b_1+a_1b_0.
\]
A field quotient uses the two equations
$b_0z_0+2b_1z_1=a_0$, $b_0z_1+b_1z_0=a_1$.
Inverses enforce their own domains. The other quotient domains
were proved above; the subtraction inverse and scalar $b_\beta$
inverse are retained explicitly.

For the logarithms, the program evaluates $G_{12}$ in blocks of
three powers and $G_3$ from its first four powers. It shares the
coefficients of
\[
 E_m(T)=\prod_{j:m_j=m}(1+TD_j^2),\qquad m=1,16,52,
\]
whose degrees are 10, 2, and 4. For a detector put
$e_r=(\mathcal A_r\delta_r)^5/F_5(t_r)$. Then
\[
 R_r=(e_r\delta_r)^{-1},\qquad
 q_1=e_r/\delta_r^3,\quad q_{16}=q_1e_r^{30},\quad
 q_{52}=q_1e_r^{102},\qquad
 W_r=E_1(q_1)E_{16}(q_{16})E_{52}(q_{52}).
\]
These identities are precisely the weights
$\mathcal A_r^{10m-5}\delta_r^{10m-8}/F_5(t_r)^{2m-1}$.
The power chain uses exponents $2,3,6,12,24,48,96$ before the
last products. The inverse of $e_r$ adds no restriction because
$\mathcal A_r,\delta_r,F_5(t_r)$ are already nonzero.
Writing $W_r=w_0+w_1\sqrt2$, the trace step is
\[
 B_r=R_r\left(\frac{w_0^2-2w_1^2}{2w_0}\right)^2.
\]
It needs only the inverse of $w_0$; over this quadratic field
$w_0\ne0$ implies $W_r\ne0$ and nonzero norm.

There are 90 initial rational coordinates and 47 structural residuals:
\[
\begin{array}{l|r|r}
 &\text{coordinates}&\text{structural residuals}\\\hline
 \text{input}&1&0\\
 \text{nine halfpoints}&36&18\\
 \text{common slope}&1&0\\
 \text{nine coordinate poles}&36&18\\
 \text{scalar pole with preparation}&5&2\\
 \text{private anchor}&3&1\\
 \text{Boolean correction}&8&8\\\hline
 \text{total}&90&47
\end{array}
\]
The complete gate list has 1214 additional coordinates and 1214
defining equations. Thus the raw system has 1304 coordinates and
1261 residuals. These are finite program counts, checked from the
full saved list, rather than estimates for an unspecified circuit.

Before elimination, apply the private affine change
\begin{equation}\label{eq:v5-anchor-translation}
 u=\mathsf S+\frac{U-3V+3W}{4},\qquad
 v=\frac{\mathsf S}{2}+\frac V4,\qquad
 w=\frac{3\mathsf S}{2}+\frac W4.
\end{equation}
Its determinant is $1/64$, with inverse
\[
 U=4u+12v-12w+8\mathsf S,\quad V=4v-2\mathsf S,\quad
 W=4w-6\mathsf S.
\]
If $T=f_0^2+f_0f_1-f_1^2$, sixteen times the old anchor after this
substitution is
\begin{equation}\label{eq:v5-translated-anchor}
 -48T^2+U^2-6UV+6UW+15V^2-18VW+7W^2+8\mathsf S U.
\end{equation}
This is the normalization used here. The three changed coordinates
occur nowhere else. Multiplying this one residual by sixteen
preserves its zero set; no identity of differently normalized raw
sums of squares is asserted.

\begin{proposition}[Projection-preserving compiler]
\label{prop:v5-compiler-projection}
The gate program, the private anchor transformation, and the
certified 32 eliminations preserve the rational input projection.
The resulting system has exactly 1272 active coordinates and
1229 quadratic residuals.
\end{proposition}
\begin{proof}
An assignment gate has a unique value once its earlier inputs are
given. Normalizing and reusing the homogeneous quadratic part
therefore changes no rational solution. The determinant of a
field quotient is $b_0^2-2b_1^2$, and its nonzero domain was
proved or explicitly imposed. Thus all gate values are the
specified arithmetic values.

Each elimination in the certificate chooses a row $aY+q=0$,
where $a\in\Q^*$ and $q$ is independent of $Y$, substitutes
$Y=-q/a$ into every other row, and removes the defining row.
It is accepted only after checking that every resulting row has
coordinate degree at most two. Insertion of $Y=-q/a$ and deletion
are inverse maps on common zero sets. The certificate performs
32 such operations, no row additions, and no omission of a free
coordinate. It retains the input and all three private anchor
coordinates. Clearing each row's rational coefficient denominators
and dividing by its integer content also preserves its zero set.
The full incidence check gives
\[
 90+1214-32=1272,\qquad47+1214-32=1229,
\]
and verifies that each retained coordinate occurs. The affine
anchor map is invertible, so it has the same projection property.
\end{proof}

Let $q_1,\ldots,q_{1229}$ be the resulting primitive integer
residuals, with the specified integer $\mathsf S$ in their
coefficients, and define
\begin{equation}\label{eq:v5-final-polynomial}
 P_{1272}(X_0,\ldots,X_{1271})=\sum_{\nu=1}^{1229}q_\nu^2.
\end{equation}
The full table is \nolinkurl{quartics/quartic-1272/polynomial.json}.
An entry with coordinate list $[i_1,\ldots,i_d]$, integer
\texttt{coefficient} $c$, and \texttt{sPower} $k$ means
\[
 c\mathsf S^kX_{i_1}\cdots X_{i_d}.
\]
Repeated indices record powers. Sum every entry; there are no
unlisted gate values in this definition. The table has SHA-256
\begin{center}\small
\nolinkurl{02e4771ba6e6c2abc749e33e8440cc256939f48005d5102fba7d9853489c8586}.
\end{center}
This is a finite specification of an ordinary integer polynomial.
The power notation in \eqref{eq:v5-fixed-constant} specifies fixed
integers exactly and adds neither variables nor polynomial degree.

\begin{proposition}[Exact table statistics]\label{prop:v5-table-statistics}
The table in \eqref{eq:v5-final-polynomial} has the degree, support,
content, and height asserted in Theorem~\ref{thm:v5-1272-polynomial}.
\end{proposition}
\begin{proof}
Independent exact convolution of all 1229 residual squares gives
48001 table entries on 48001 different coordinate monomials.
Of these, 47993 have $\mathsf S$-independent integer coefficients,
seven have coefficients $c\mathsf S$ with $0<|c|\le768$, and one
has coefficient $64\mathsf S^2$. No monomial receives different
powers of $\mathsf S$, so none disappears on specialization.
The largest absolute $\mathsf S$-independent coefficient has bit
length 5339, and is therefore less than $2^{5339}$.

The three private anchor coordinates remain private through the
eliminations. In the final numbering $U=X_{79}$. Squaring
\eqref{eq:v5-translated-anchor} shows that the coefficient of
$X_{79}^2$ is exactly $64\mathsf S^2$; it is the only coefficient
with power two in $\mathsf S$. Since
$\mathsf S=3^{2^{27}}>2^{2^{27}}$, this exceeds both
$2^{5339}$ and $768\mathsf S$. Thus it is the exact height.
The coefficient of $X_0^4$ is one, proving both degree four and
integer coefficient content one. Finally
\[
 \log_2(64\mathsf S^2)<6+\tfrac85\,2^{28}<429496736,
\]
using $3^5<2^8$.
\end{proof}

The saved construction contains the raw source, all 1214 gate
definitions with dependency order, the 32-step elimination
certificate, row normalizations, the compact coordinate map,
and the final residual and coefficient tables. The independent
checker imports neither producer nor compiler code. It checks
each gate identity, replays every elimination, checks every
coefficient by a separate convolution, and verifies the fixed
integer support and height bounds. A fresh isolated reconstruction
reproduces the displayed table byte for byte.

The source review is a separate part of the argument. It matches
the implemented formulas and guards to the source just proved.
Two separately written evaluators, including direct sixteen-factor
detector products, provide additional exact finite-field and
rational checks. Those sampled checks are not substituted for
the domain proofs or the arithmetic theorem.

\begin{proof}[Proof of Theorem~\ref{thm:v5-1272-polynomial}]
The finite assertions follow from Propositions
\ref{prop:v5-compiler-projection} and \ref{prop:v5-table-statistics}.
Over $\Q$, the sum in \eqref{eq:v5-final-polynomial} is zero
exactly when all its residuals are zero. Its rational input
projection is therefore the one in
Theorem~\ref{thm:v5-source-projection}, namely $\Z$.
\end{proof}
